\chapter{Ingrown Toenails (Onychocryptosis)}

\section*{Definition}
An ingrown toenail is a condition where the nail grows into the surrounding skin, causing pain, redness, swelling, and often infection. It most commonly affects the great toe.

\section*{What Happens in Your Body}
The nail plate pierces the lateral nail fold, causing a foreign-body inflammatory response. This leads to pain, edema, and granulation tissue formation. If the skin barrier is compromised, secondary bacterial infection can occur. Improper nail trimming, tight footwear, and genetic factors contribute to ingrowth.

\section*{Causes}
\begin{itemize}
\item Improper nail trimming (rounded corners instead of straight across)
\item Tight shoes that compress the toes
\item Trauma: stubbing the toe, dropping objects on the foot
\item Genetic predisposition: wide nail fold or curved nails
\item Poor foot hygiene
\item Excessive sweating
\item Obesity (increased pressure on toes)
\item Repetitive activities (kicking, running)
\end{itemize}

\section*{When to See a Doctor}
\begin{itemize}
\item Pain along the edge of the toenail
\item Redness and swelling of the nail fold
\item Tenderness to touch
\item Pus or drainage from the nail edge
\item Overgrowth of skin (granulation tissue) around the nail
\item Difficulty wearing shoes
\item Chronic or recurrent episodes
\end{itemize}

\section*{Related Symptoms}
\begin{itemize}
\item Paronychia (infection of the nail fold)
\item Cellulitis of the toe
\item Osteomyelitis (bone infection, rare)
\item Fungal nail infection
\item Subungual hematoma
\end{itemize}

\section*{Self-Care / Home Management}
\begin{itemize}
\item Soak the foot in warm, soapy water 3-4 times daily
\item Gently lift the nail edge and place a small piece of cotton or dental floss under it
\item Apply antibiotic ointment
\item Wear open-toed shoes or sandals
\item Take OTC pain relievers
\item Do NOT cut the nail at the corner or dig into the nail fold
\item Trim nails straight across
\end{itemize}

\section*{How Your Doctor Will Evaluate You}
\begin{itemize}
\item Clinical examination of the toe and nail
\item Assess for signs of infection (erythema, purulence, warmth)
\item X-ray if osteomyelitis suspected in severe cases
\end{itemize}

\section*{Treatment Options}
\begin{itemize}
\item Conservative: warm soaks, cotton wick placement, antibiotics
\item Partial nail avulsion (removal of the ingrown nail edge)
\item Matrixectomy (chemical ablation with phenol) to prevent recurrence
\item Complete nail avulsion for recurrent or severe cases
\item Oral antibiotics for cellulitis
\item Surgical correction by a podiatrist
\end{itemize}

\section*{References}

\begin{thebibliography}{99}
\bibitem{uptodate-ingrown} Eekhof JAH, et al. Ingrown toenail. In: UpToDate, Post TW (Ed), UpToDate, Waltham, MA; 2025.
\bibitem{coughlin} Coughlin MJ, et al. \emph{Mann's Surgery of the Foot and Ankle}. 10th ed. Elsevier; 2023. Chapter 29: Ingrown Toenails.
\bibitem{heidelbaugh} Heidelbaugh JJ, et al. Management of ingrown toenails. \emph{Am Fam Physician}. 2023;107(3):275-283.
\bibitem{declerck} Declerck G, et al. Ingrown toenails: a systematic review of treatment. \emph{J Foot Ankle Surg}. 2024;63(1):98-107.
\bibitem{ekhof} Eekhof JAH, et al. Conservative versus surgical treatment of ingrown toenails. \emph{Br J Gen Pract}. 2022;72(721):380-386.
\bibitem{rozeboom} Rozeboom AM, et al. Phenol matrixectomy for ingrown toenail. \emph{Dermatol Surg}. 2023;49(6):567-574.
\end{thebibliography}
